\setcounter{ex}{0}

% --- CÂU 1 ---
\begin{ex}
	Cho $\sin\alpha = \dfrac{5}{13}$ với $\dfrac{\pi}{2} < \alpha < \pi$. Giá trị của $\cos 2\alpha$ bằng
	\choice
	{$\dfrac{144}{169}$}
	{$\dfrac{120}{169}$}
	{\True $\dfrac{119}{169}$}
	{$-\dfrac{119}{169}$}
	\loigiai{
		Áp dụng công thức nhân đôi $\cos 2\alpha = 1 - 2\sin^2\alpha$, ta có:
		$$\cos 2\alpha = 1 - 2 \cdot \left(\dfrac{5}{13}\right)^2 = 1 - \dfrac{50}{169} = \dfrac{119}{169}.$$
	}
\end{ex}

% --- CÂU 2 ---
\begin{ex}
	Cho $\sin a = \dfrac{4}{5}$, $\sin b = \dfrac{5}{13}$ với $0 < a, b < \dfrac{\pi}{2}$. Giá trị của $\cos(a - b)$ bằng
	\choice
	{\True $\dfrac{56}{65}$}
	{$\dfrac{33}{65}$}
	{$\dfrac{16}{65}$}
	{$\dfrac{63}{65}$}
	\loigiai{
		Vì $0 < a, b < \dfrac{\pi}{2}$ nên $\cos a > 0$ và $\cos b > 0$. Ta có:
		$$\cos a = \sqrt{1 - \sin^2 a} = \dfrac{3}{5}; \quad \cos b = \sqrt{1 - \sin^2 b} = \dfrac{12}{13}.$$
		Khi đó:
		$$\cos(a - b) = \cos a\cos b + \sin a\sin b = \dfrac{3}{5} \cdot \dfrac{12}{13} + \dfrac{4}{5} \cdot \dfrac{5}{13} = \dfrac{36 + 20}{65} = \dfrac{56}{65}.$$
	}
\end{ex}

% --- CÂU 3 ---
\begin{ex}
	Tập xác định của hàm số $y = \dfrac{1}{1 - \cos 2x}$ là
	\choice
	{$\mathbb{R} \setminus \left\{\dfrac{\pi}{2} + k\pi,\ k \in \mathbb{Z}\right\}$}
	{$\mathbb{R} \setminus \{k2\pi,\ k \in \mathbb{Z}\}$}
	{$\mathbb{R} \setminus \left\{\dfrac{k\pi}{2},\ k \in \mathbb{Z}\right\}$}
	{\True $\mathbb{R} \setminus \{k\pi,\ k \in \mathbb{Z}\}$}
	\loigiai{
		Hàm số xác định khi và chỉ khi:
		$$1 - \cos 2x \ne 0 \iff \cos 2x \ne 1 \iff 2x \ne k2\pi \iff x \ne k\pi\ (k \in \mathbb{Z}).$$
		Vậy tập xác định của hàm số là $\mathscr{D} = \mathbb{R} \setminus \{k\pi,\ k \in \mathbb{Z}\}$.
	}
\end{ex}

% --- CÂU 4 ---
\begin{ex}
	Chu kì tuần hoàn của hàm số $y = \sin 3x$ là
	\choice
	{$6\pi$}
	{\True $\dfrac{2\pi}{3}$}
	{$2\pi$}
	{$\dfrac{\pi}{3}$}
	\loigiai{
		Hàm số $y = \sin 3x$ tuần hoàn với chu kì $T = \dfrac{2\pi}{3}$.
	}
\end{ex}

% --- CÂU 5 ---
\begin{ex}
	Nghiệm của phương trình $\sqrt{3}\tan x + 3 = 0$ là
	\choice
	{$x = \dfrac{\pi}{3} + k\pi\ (k \in \mathbb{Z})$}
	{\True $x = -\dfrac{\pi}{3} + k\pi\ (k \in \mathbb{Z})$}
	{$x = -\dfrac{\pi}{6} + k\pi\ (k \in \mathbb{Z})$}
	{$x = \dfrac{\pi}{6} + k\pi\ (k \in \mathbb{Z})$}
	\loigiai{
		Ta có:
		$$\sqrt{3}\tan x + 3 = 0 \iff \tan x = -\sqrt{3} \iff x = -\dfrac{\pi}{3} + k\pi\ (k \in \mathbb{Z}).$$
	}
\end{ex}

% --- CÂU 6 ---
\begin{ex}
	Số nghiệm của phương trình $\sin x = -\dfrac{1}{2}$ trên khoảng $(0; 4\pi)$ là
	\choice
	{\True $4$}
	{$3$}
	{$2$}
	{$5$}
	\loigiai{
		Ta có:
		$$\sin x = -\dfrac{1}{2} \iff \hoac{x &= -\dfrac{\pi}{6} + k2\pi \\ x &= \dfrac{7\pi}{6} + k2\pi}\ (k \in \mathbb{Z}).$$
		Xét các nghiệm thuộc khoảng $(0; 4\pi)$:
		\begin{itemize}
			\item Với $x = -\dfrac{\pi}{6} + k2\pi \in (0; 4\pi) \implies k \in \{1; 2\}$ (có $2$ nghiệm là $\dfrac{11\pi}{6}; \dfrac{23\pi}{6}$).
			\item Với $x = \dfrac{7\pi}{6} + k2\pi \in (0; 4\pi) \implies k \in \{0; 1\}$ (có $2$ nghiệm là $\dfrac{7\pi}{6}; \dfrac{19\pi}{6}$).
		\end{itemize}
		Vậy phương trình có $4$ nghiệm trên khoảng $(0; 4\pi)$.
	}
\end{ex}

% --- CÂU 7 ---
\begin{ex}
	Giá trị lớn nhất của hàm số $y = 5 - 3\cos\left(2x - \dfrac{\pi}{3}\right)$ bằng
	\choice
	{$5$}
	{$3$}
	{\True $8$}
	{$2$}
	\loigiai{
		Vì $-1 \le \cos\left(2x - \dfrac{\pi}{3}\right) \le 1, \forall x \in \mathbb{R}$ nên:
		$$2 \le 5 - 3\cos\left(2x - \dfrac{\pi}{3}\right) \le 8.$$
		Vậy giá trị lớn nhất của hàm số bằng $8$ khi $\cos\left(2x - \dfrac{\pi}{3}\right) = -1$.
	}
\end{ex}

% --- CÂU 8 ---
\begin{ex}
	Cho cấp số cộng $(u_n)$ có $u_3 = 7$ và $u_8 = -3$. Công sai $d$ của cấp số cộng là
	\choice
	{\True $-2$}
	{$-3$}
	{$2$}
	{$3$}
	\loigiai{
		Ta có $u_8 = u_3 + 5d \iff -3 = 7 + 5d \iff 5d = -10 \iff d = -2$.
	}
\end{ex}

% --- CÂU 9 ---
\begin{ex}
	Cho cấp số cộng $(u_n)$ có $u_1 = 2$ và công sai $d = 3$. Số $152$ là số hạng thứ mấy của cấp số cộng?
	\choice
	{\True $51$}
	{$49$}
	{$50$}
	{$52$}
	\loigiai{
		Giả sử $152$ là số hạng thứ $n$ của cấp số cộng ($n \in \mathbb{N}^*$). Ta có:
		$$u_n = u_1 + (n - 1)d = 152 \iff 2 + 3(n - 1) = 152 \iff n - 1 = 50 \iff n = 51.$$
		Vậy số $152$ là số hạng thứ $51$ của cấp số cộng.
	}
\end{ex}

% --- CÂU 10 ---
\begin{ex}
	Cho hình chóp $S.ABCD$ có đáy $ABCD$ là hình thang ($AB \parallel CD,\ AB = 2CD$), $E = AD \cap BC$, $M, N$ lần lượt là trung điểm của $SD, SC$ (hình vẽ dùng chung cho câu 10, 11, 12). Giao tuyến của hai mặt phẳng $(SAD)$ và $(SBC)$ là
	\begin{center}
		\begin{tikzpicture}[scale=0.85, font=\footnotesize, line join=round, line cap=round]
			\coordinate (A) at (0,0);
			\coordinate (B) at (4,0);
			\coordinate (D) at (1,1.1);
			\coordinate (C) at (3,1.1);
			\coordinate (E) at (2,2.2);
			\coordinate (O) at (intersection of A--C and B--D);
			\coordinate (S) at (2,3.6);
			\coordinate (M) at ($(S)!0.5!(D)$);
			\coordinate (N) at ($(S)!0.5!(C)$);
			\draw[thick] (S)--(A) (S)--(B) (A)--(B);
			\draw[dashed] (S)--(D) (S)--(C) (A)--(D)--(C)--(B) (A)--(C) (B)--(D) (D)--(E)--(C) (S)--(E) (M)--(N);
			\foreach \p/\pos in {S/above, A/below left, B/below right, C/right, D/left, O/below, M/above left, N/above right, E/below}
				\node[\pos] at (\p) {$\p$};
			\foreach \p in {S,A,B,C,D,O,M,N,E}
				\fill (\p) circle (1.2pt);
		\end{tikzpicture}
	\end{center}
	\choice
	{\True $SE$}
	{$SO$}
	{$SA$}
	{$SC$}
	\loigiai{
		Ta có $S \in (SAD) \cap (SBC)$ và $E = AD \cap BC \implies E \in (SAD) \cap (SBC)$.\\
		Do đó giao tuyến của hai mặt phẳng $(SAD)$ và $(SBC)$ là đường thẳng $SE$.
	}
\end{ex}

% --- CÂU 11 ---
\begin{ex}
	Gọi $F$ là giao điểm của đường thẳng $SA$ và mặt phẳng $(BCM)$. Tỉ số $\dfrac{SF}{SA}$ bằng
	\choice
	{$\dfrac{2}{3}$}
	{\True $\dfrac{1}{3}$}
	{$\dfrac{1}{2}$}
	{$\dfrac{1}{4}$}
	\loigiai{
		Vì $E = AD \cap BC$ nên $E \in BC \subset (BCM)$, suy ra $(BCM) \cap (SAD) = EM$.\\
		Do đó giao điểm của $SA$ và $(BCM)$ là $F = SA \cap EM$.\\
		Trong tam giác $ABE$, vì $CD \parallel AB$ và $CD = \dfrac{1}{2}AB$ nên $D$ là trung điểm của $AE$.\\
		Trong tam giác $SAE$, kẻ $DK \parallel EF\ (K \in SA)$:
		\begin{itemize}
			\item Xét tam giác $AEF$ có $D$ là trung điểm của $AE$ và $DK \parallel EF$ nên $K$ là trung điểm của $AF \implies AF = FK$.
			\item Xét tam giác $SDK$ có $M$ là trung điểm của $SD$ và $MF \parallel DK$ nên $F$ là trung điểm của $SK \implies SF = FK$.
		\end{itemize}
		Suy ra $SF = FK = KA \implies \dfrac{SF}{SA} = \dfrac{1}{3}$.
	}
\end{ex}

% --- CÂU 12 ---
\begin{ex}
	Khẳng định nào sau đây đúng?
	\choice
	{$MN \parallel (SBC)$}
	{$MN \parallel (SAD)$}
	{$MN \parallel (SCD)$}
	{\True $MN \parallel (SAB)$}
	\loigiai{
		Vì $M, N$ lần lượt là trung điểm của $SD, SC$ nên $MN$ là đường trung bình của tam giác $SCD \implies MN \parallel CD$.\\
		Mà $CD \parallel AB$ nên $MN \parallel AB$. Do $AB \subset (SAB)$ và $MN \not\subset (SAB)$ nên $MN \parallel (SAB)$.
	}
\end{ex}
